\documentclass{article}
\usepackage{amsmath, amssymb}
\usepackage{amsthm}

\author{Michael Naumov}
\date{}
\setlength{\parindent}{0pt}
\setlength{\parskip}{0.5em}


\title{1.2.4-48} % chktex 8

\begin{document}
\maketitle

\subsubsection*{(a)}

\paragraph*{1. \$\{\} n < 0 \{\}\$}

Counterexample: \( m = 0, n = -1 \)

\( \lfloor{\frac{m+n-1}{n}}\rfloor = \lfloor{\frac{0-1-1}{-1}}\rfloor = 2 \)

\( \lceil{\frac{m}{n}}\rceil = \lceil{\frac{0}{-1}}\rceil = 0 \)

\paragraph*{2. \$\{\} n = 1 \{\}\$}

\( \lfloor{\frac{m+n-1}{n}}\rfloor = \lfloor{\frac{m}{1}}\rfloor = m \)

\( \lceil{\frac{m}{n}}\rceil = \lceil{\frac{m}{1}}\rceil = m \)

\paragraph*{3. \$\{\} n \textbackslash\{\}geq 2 \{\}\$}

Let \( m = qn+r, 0 \leq r \leq n - 1 \)

\( \lfloor{\frac{m+n-1}{n}}\rfloor = \lfloor{\frac{qn+r+n-1}{n}}\rfloor = q + 1 + \lfloor{\frac{r-1}{n}}\rfloor \)

##### 3.1.\ \( r = 0 \)

\( \lfloor{\frac{m+n-1}{n}}\rfloor = q + 1 + \lfloor{-\frac{1}{n}}\rfloor = q + 1 - 1 = q \)

\( \lceil{\frac{m}{n}}\rceil = \lceil{\frac{qn}{n}}\rceil = q \)

##### 3.2. \( 1 \leq r \leq n - 1 \)

\( \lfloor{\frac{m+n-1}{n}}\rfloor = q + 1 + \lfloor{\frac{r-1}{n}}\rfloor = q + 1 \)

\( \lceil{\frac{m}{n}}\rceil = \lceil{\frac{qn+r}{n}}\rceil = q + \lceil{\frac{r}{n}}\rceil = q + 1 \)

\subsubsection*{(b)}

\( \lfloor{\frac{n+2 - \lfloor{\frac{n}{25}}\rfloor}{3}}\rfloor = \lceil{\frac{n - \lfloor{\frac{n}{25}}\rfloor}{3}}\rceil = \lceil{\frac{n + \lceil{-\frac{n}{25}}\rceil}{3}}\rfloor = \lceil{\frac{n + \lceil{-\frac{n}{25}}\rceil}{3}}\rceil = \lceil{\frac{\lceil{n-\frac{n}{25}}\rceil}{3}}\rceil = \lceil{\frac{\lceil{\frac{24n}{25}}\rceil}{3}}\rceil = \lceil{\frac{8n}{25}}\rceil = \lfloor{\frac{8n+24}{25}}\rfloor \)


\end{document}
